\section{NN-S014: a logarithmic cancellation bound and a uniform distinguisher}

\paragraph{Statement and exact inputs.}
There is an absolute constant $c>0$ such that every monotone map
$S:\{0,1\}^n\to\{0,1\}^{n+1}$, for all sufficiently large $n$,
has a test depending on at most two output coordinates with distinguishing
advantage at least $c\log n/n^2$ against the uniform distribution.
The test is a coordinate or the AND of two coordinates, with its
acceptance convention possibly complemented. For a uniformly
polynomial-time computable family $S$, a uniform randomized polynomial-time
algorithm selects such a test to within an additive $O(n^{-3})$ loss,
with failure probability at most $2(n+2)^2e^{-2n^2}$.

Consequently, for a uniformly constructible polynomial-size $G(R)$ circuit
family over any fixed finite nonzero commutative ring, with Boolean
amplitude output and positive stretch, ordinary security against
\emph{uniform randomized polynomial-time} distinguishers implies that
some final output equals the supported zero on
\[
                 \Omega(2^n\log n/n^2)
\]
seed inputs. The same bound holds for nonuniform-secure generators without
the uniform circuit-construction hypothesis. This sharpens the quantitative
bound and removes the nonuniform-adversary restriction of NN-S013 for
uniform circuit families. It is a necessary condition for pseudorandomness,
not an existence assertion or a circuit-size lower bound.

The classical inputs are Goldreich--Izsak, \emph{Monotone Circuits:
One-Way Functions versus Pseudorandom Generators}, Theory of Computing
8 (2012), 231--238: Corollary~4.4, pp.~235--236, and Theorem~4.5,
p.~236, followed by the proof of Theorem~1.2 in Section~4.2,
pp.~236--237. The logarithm is already present in their covariance
estimate; the following proof changes the marginal-bias threshold.
Thus no priority claim is made for a new classical correlation theorem.
The new split-circuit and uniform-selection consequences are proved below.
The primary source is
\url{https://theoryofcomputing.org/articles/v008a010/v008a010.pdf}.

\paragraph{Proof of the sharper two-coordinate estimate.}
Write $f_i$ for the $i$th coordinate of $S$ and $I_i(j)$ for its
influence at input coordinate $j$. The KKL estimate in the cited
Corollary~4.4 gives an absolute $a>0$ such that every function with bias
$o(\log n/n)$ has some influence at least $a\log n/n$, for all
sufficiently large $n$. The constant can be chosen uniformly here:
changing a function on a fraction $\delta$ of inputs makes it exactly
balanced and changes each influence by at most $2\delta$, as proved
immediately before the cited corollary.

Suppose that every marginal bias is at most
$\delta=c_0\log n/n^2$, where the absolute $c_0>0$ will be fixed below.
Choose one such influential input for each of the $n+1$ outputs.
Two distinct outputs, say $f_i,f_k$, share a chosen input. The
Talagrand inequality in the cited Theorem~4.5 gives
\[
 \operatorname{Cov}(f_i,f_k)
 \ge c_T\frac{T}{\log(e/T)},\qquad
 T=\sum_j I_i(j)I_k(j)
 \ge a^2\frac{(\log n)^2}{n^2}.
\]
The expression on the right is increasing for $0<T\le1$ and is
at least $c_1\log n/n^2$ for an absolute $c_1>0$ and all large $n$.
For completeness, $T\le1$: for a monotone Boolean function in the
uniform cube, the absolute value of its first Fourier coefficient is one half of its
influence, so Parseval gives $\sum_j I_i(j)^2\le4\operatorname{Var}(f_i)
\le1$; Cauchy--Schwarz proves the bound for $T$.
The displayed lower estimate also follows directly by substituting
$a^2(\log n)^2/n^2$ and bounding its logarithmic denominator by
$3\log n$ for sufficiently large $n$.

Since both marginal means are at least $1/2-\delta$,
\[
 \mathbb E(f_if_k)-\frac14
 \ge c_1\frac{\log n}{n^2}-\delta+\delta^2.
\]
Fix $c_0=c_1/4$. If a marginal has bias greater than $\delta$,
that coordinate gives advantage greater than $\delta$. Otherwise
the AND of the two coordinates above has advantage at least
$3c_1\log n/(4n^2)$. Decreasing the universal constant if needed
proves the assertion for every sufficiently large $n$. The proof
uses no computation of the influential coordinates.

\paragraph{Uniform selection of the test.}
Let $\mathcal T$ consist of the $n+1$ coordinate tests and all
two-coordinate AND tests. Its size is at most $(n+2)^2$.
For each $D\in\mathcal T$, its uniform mean $u_D$ is known exactly:
it is $1/2$ or $1/4$. Generate $L=n^8$ independent samples
$S(U_n)$ and compute every empirical mean $\widehat\mu_D$.
Put $\varepsilon=n^{-3}$. The elementary Bernoulli exponential
bound and a union bound give
\[
 \Pr\left[\max_{D\in\mathcal T}
       |\widehat\mu_D-\mathbb E D(S(U_n))|>\varepsilon\right]
 \le 2(n+2)^2e^{-2n^2}=:\eta_n.                 \tag{S014.1}
\]
Here the Bernoulli bound follows by exponential Markov applied to
independent centred samples and the inequality
$\mathbb E e^{s(X-\mathbb EX)}\le e^{s^2/8}$ for $0\le X\le1$.
One proof of this latter inequality differentiates the logarithmic
moment-generating function twice: its second derivative is the variance
under an exponential tilt and is at most $1/4$; its value and first
derivative vanish at zero. Optimizing $s$ gives each tail
$e^{-2L\varepsilon^2}$, so no additional probabilistic premise is hidden.

Choose a test maximizing $|\widehat\mu_D-u_D|$. Use $D$ itself
when the empirical bias is nonnegative and $1-D$ otherwise.
On the event complementary to (S014.1), the chosen, oriented test
$D_*$ satisfies
\[
 \mathbb E D_*(S(U_n))-\mathbb E D_*(U_{n+1})
 \ge c\frac{\log n}{n^2}-2\varepsilon.           \tag{S014.2}
\]
Indeed a genuinely best test loses at most $\varepsilon$ on
estimation, empirical maximization cannot decrease that value,
and the chosen test loses at most a further $\varepsilon$.
The training samples are independent of the challenge on which
the test is evaluated. The entire selection and evaluation algorithm
is uniform and polynomial time when $S$ is uniform and polynomial time.
Neither a hard-coded influential coordinate nor nonuniform advice is
required. The exponent eight is only a convenient fixed polynomial;
no optimal sample complexity is asserted.

\paragraph{Transfer to supported zeros.}
Let $F$ be the first $n+1$ Boolean amplitude outputs and $S$ their
support outputs. As in NN-S013, $S$ is monotone and
$F_i(x)\ne S_i(x)$ holds exactly on
$B_i=\{x:C_i(x)=e\}$. Put
$b_n=\max_{i\le n+1}\Pr[U_n\in B_i]$.
Every selected test depends on at most two coordinates. For each fixed
training outcome, coupling $F$ and $S$ on the same seed changes its
expectation by at most $2b_n$; complementing the test does not change
this bound. Average (S014.2) over the independent training samples.
On the exceptional event any bias is at least $-1$. Therefore the
uniform randomized distinguisher satisfies the unconditional inequality
\[
 \mathbb E D_*(F(U_n))-\mathbb E D_*(U_{n+1})
 \ge c\frac{\log n}{n^2}-2b_n-2n^{-3}-2\eta_n.  \tag{S014.3}
\]
If $F$ is secure against uniform randomized polynomial time, the
absolute value of the left side is negligible. Thus
\[
 b_n\ge\frac{c\log n}{2n^2}-n^{-3}-\eta_n
                   -\operatorname{negl}(n)
      =\Omega(\log n/n^2).
\]
Multiplying by $2^n$ gives the cardinality claim. For nonuniform security,
one may simply select the test from the first part as advice, obtaining
the same conclusion without the training step. Truncation preserves
pseudorandomness, so every polynomial output length greater than $n$
is covered. No use of E, A, or M occurs in this distributional theorem.
\hfill$\square$

\paragraph{Geometric and complexity dictionary.}
For each input length, the Boolean input algebra is the coordinate
algebra of the finite split scheme with $2^n$ points. Each $B_i$ is a
clopen subscheme, and its finite free rank is $\#B_i$.
The conclusion is therefore a lower bound on the rank of at least one
supported-zero output fibre, with probability measured by rank divided
by $2^n$. The finite-algebra trace convention and its base-change
identity are EGA~IV$_4$, Section~18.2.1, especially equation~(18.2.1.2),
and Section~18.2.2, printed p.~111. This is an explicit finite-scheme
dictionary; no l-adic weight or purity theorem is used to obtain the
probability bound.

The construction of the clopen fibre is implicit in the circuit;
its exponentially many points need not be enumerated. The theorem
improves a necessary cancellation condition and supplies a uniform
attack when that condition fails. It gives neither a lower bound on
the total number of circuit gates nor a proof of P versus NP or
VP versus VNP. E's separate subexponential inversion consequences
retain their original sparsity and locality hypotheses.
